\section{Koide and Local Kernel Construction}
\label{v480:koide-kernel}

This section supplies explicit arguments for the two proposed compression
sites. It supersedes the claim that cyclic symmetry alone proves the
inherited Koide theorem. It also constructs a relative kernel model,
without declaring the full prospective native zero-kernel layer active.

\subsection{The Exact Additional Koide Condition}
\label{v480:koide-condition}

Let $x_i=\sqrt{m_i}\geq0$, $x\ne0$, and write
\[
 A=\frac{x_0+x_1+x_2}{3}>0,\qquad
 u=A(1,1,1),\quad d=x-u,\quad \sum_i d_i=0.
\]
The cyclic permutation preserves the one-dimensional symmetric line
and its two-dimensional orthogonal plane. It preserves their norms
separately; it does not prescribe their ratio.

\begin{proposition}[Koide as an equal-norm condition]
\label{v480:koide-equal-norm}
For these nonnegative principal roots,
\[
 K=\frac{\sum_i x_i^2}{(\sum_i x_i)^2}
   =\frac13+\frac{\|d\|^2}{9A^2}.
\]
Consequently $K=2/3$ if and only if
$\|d\|^2=\|u\|^2=3A^2$.
\end{proposition}
\begin{proof}
Since $u\perp d$, $\|x\|^2=3A^2+\|d\|^2$, while
$(\sum_i x_i)^2=9A^2$. Division proves both assertions.
\end{proof}

Thus Cauchy--Schwarz supplies only $K\geq1/3$. For example,
the cyclic orbit of $x=(1,2,3)$ preserves both sums used in the
inherited argument, but gives $K=14/36=7/18$, not $2/3$.
Even the fully symmetric choice $x=(1,1,1)$ gives $K=1/3$.
The missing model obligation is an independently derived equal-norm
law, not a normalization of the Cauchy--Schwarz inequality by a
count of sectors.

\begin{proposition}[A sufficient HC-shaped mass-root realization]
\label{v480:koide-hc-mixing}
In a declared real three-root representation, let
$e=(1,1,1)/\sqrt3$, let $J_m$ be an orthogonal involution, and put
$C_m=(I+J_m)/2$. Suppose $x=\alpha C_m e$ is nonzero and
componentwise nonnegative, with $\alpha>0$. For
$h=\langle e,J_m e\rangle$, one has $h>-1$ and
\[
 K=\frac{2}{3(1+h)}.
\]
Hence Koide holds precisely when $\langle e,J_m e\rangle=0$.
\end{proposition}
\begin{proof}
An orthogonal involution is self-adjoint, so $C_m$ is an
orthogonal projector. Thus
$\|C_m e\|^2=\langle e,C_m e\rangle=(1+h)/2>0$ and
$(\sum_i x_i)^2=3\alpha^2(1+h)^2/4$. Dividing
$\|x\|^2=\alpha^2(1+h)/2$ by this denominator gives the formula.
\end{proof}

This condition is realizable: take $v=(2,-1,-1)/\sqrt6$, a unit
vector orthogonal to $e$ with $e+v$ positive. Let $J_m$ exchange $e$ and $v$, and let it
act as $+1$ on their orthogonal complement. Then
$C_m e=(e+v)/2$ has equal symmetric and orthogonal norms.
The direction of $v$ still determines a free phase. Conversely,
$J_m=I$ is an equally valid orthogonal involution but gives $K=1/3$.
Therefore HC projection axioms alone do not force Koide. The
representation of the native HC on mass roots, positivity and the
orthogonality condition must be derived rather than selected to fit.

\begin{proposition}[Radius and phase in the three-root parametrization]
\label{v480:koide-radius-phase}
Every nonzero nonnegative three-root vector admits
\[
 x_k=A\bigl(1+r\cos(\theta+2\pi k/3)\bigr),\qquad r\geq0,
\]
with $K=1/3+r^2/6$. The condition $K=2/3$ is precisely $r=\sqrt2$;
it does not fix $A$ or $\theta$.
\end{proposition}
\begin{proof}
The vectors $(\cos(2\pi k/3))_{k=0}^2$ and
$(\sin(2\pi k/3))_{k=0}^2$ form an orthogonal basis of the
zero-sum plane, each with squared norm $3/2$. Expressing $d/A$
in that basis gives coefficients $r\cos\theta$ and
$-r\sin\theta$. Therefore $\|d\|^2=(3/2)A^2r^2$ and the
previous proposition gives the formula for $K$. If $r=0$ the
phase is arbitrary.
\end{proof}

The root formula must retain $1+r\cos(\theta+2\pi k/3)\geq0$.
Squaring a negative bracket and then calling it a principal square
root changes the denominator to a sum of absolute values; the
displayed identity for the physical $K$ no longer follows.
The phase ansatz $\theta=\pi-\sqrt2\arcsin\beta$ also does not
give mass degeneracy at $\beta=0$: with $r=\sqrt2$ it gives
\[
 (m_0,m_1,m_2)/A^2
 =\bigl(3-2\sqrt2,\;3/2+\sqrt2,\;3/2+\sqrt2\bigr),
\]
and the first signed-root bracket is negative. A degenerate
nonzero three-mass limit has $r=0$ and $K=1/3$, so it cannot
be obtained by changing only the phase at fixed Koide radius.

\begin{proposition}[Cyclic-equivariant linear maps preserve symmetric inputs]
\label{v480:koide-equivariant-obstruction}
Let $P$ cyclically permute three real coordinates. If a linear map
$T$ satisfies $TP=PT$, then $T(A(1,1,1))$ remains on the symmetric
line. It cannot produce a nondegenerate root hierarchy from that input.
\end{proposition}
\begin{proof}
$Pu=u$ implies $P(Tu)=T(Pu)=Tu$. The fixed space of $P$ is
exactly the line spanned by $(1,1,1)$.
\end{proof}

A hierarchy mechanism must therefore identify an asymmetric source,
a symmetry-breaking selection or a non-equivariant coupling; cyclic
symmetry alone is insufficient. A derivation of the phase law and the
equal-norm condition remains open. Using only the fixed mass inputs
quoted in the inherited text, the phase diagnostic is about
$132.73279^\circ$, while the stated $\beta=0.5493$ ansatz gives
$132.87981^\circ$. With $A$ fixed by the quoted roots, the respective
mass errors are approximately $-12.80\%$, $+1.20\%$, $-0.0666\%$.
These reproduce a historical-input comparison, not a new empirical
measurement or a physical derivation. The inherited
\texttt{thm:koide-triolic} and dependent phase-prediction wording
must be read with these limitations.

\subsection{A Quotient Kernel on a Declared Fragment}
\label{v480:kernel-fragment}

Fix a nonempty \emph{set} $A_F$ of admitted payloads for one declared
fragment, with tagged embeddings $j_t(a)=\langle t,a\rangle$ and
payload projection $p_t(j_t(a))=a$. Here $t=0$ denotes the retained
pre-shift display and $t=1$ its inhabited shifted state display.
This local notation does not undo the later shifted normal form.
Assume a supplied equivalence relation $E_F$ on $A_F$ expressing
the restricted closure-class schema. Its native interpretation is
input data, not something derived from the tag.

\begin{proposition}[Choice-free semantic quotient]
\label{v480:kernel-quotient}
The set $Q_F=A_F/E_F$ and map $q_F(a)=[a]_{E_F}$ are defined in ZF.
They satisfy $q_F(a)=q_F(b)$ if and only if $aE_Fb$.
If a total closure-relevant readout $f:A_F\to Y$ is constant on
$E_F$-classes, there is a unique $\bar f:Q_F\to Y$ with
$f=\bar f\circ q_F$.
\end{proposition}
\begin{proof}
Each class is a subset of $A_F$ defined by Separation, and their
collection is a set by Replacement. Reflexivity, symmetry and
transitivity give the equality criterion for classes. Define
\[
 \operatorname{graph}(\bar f)
 =\{\,([a]_{E_F},f(a)):a\in A_F\,\}.
\]
Class constancy makes this a function, and surjectivity of $q_F$
gives uniqueness. No representative or choice function is required.
\end{proof}

This quotient packages the supplied closure relation; it does not
compute it, prove a native closure claim or guarantee a cheaper
algorithm. The quotient kernel and the raw payload have distinct
roles: $j_t\circ p_t$ reconstructs the syntactic state, whereas
$q_F$ may identify different payloads and cannot reconstruct each
original state by itself.

\begin{proposition}[Finite representative normal forms]
\label{v480:kernel-finite-normal-form}
If $A_F$ is finite with a declared injective integer coding, let
$N_F(a)$ be the least-coded member of $[a]_{E_F}$. Then
$N_F:A_F\to A_F$ is total and idempotent, and
\[
 N_F(a)=N_F(b)\quad\Longleftrightarrow\quad aE_Fb.
\]
Every total class-constant readout factors uniquely through
$N_F(A_F)$. An admissible transport pair $aE_FT(a)$ gives
$N_F(T(a))=N_F(a)$ wherever the declared fragment contains both sides.
\end{proposition}
\begin{proof}
Each nonempty finite class has a unique least code. Equivalent
payloads have the same class and minimum. Conversely, equal minima
belong to both classes, which are then equal. The minimum of a
minimum's class is unchanged. Apply the quotient factorization
after identifying a class with its minimum; the transport assertion
uses the equality criterion.
\end{proof}

Thus total payload extraction and finite normal-form construction
can genuinely be discharged on this relative fragment. Identifying
$A_F$ with the paper's intended native kernel class still requires
the corresponding admissibility proof. Moreover, each aspect map
required to survive normalization must be constant on $E_F$-classes.
A signature distinguishing equivalent payloads does not descend;
calling it an invariant does not repair this conflict.

\subsection{Readout Domains and Routed Aggregation}
\label{v480:kernel-factor-and-route}

The definition of closure-compatible readout constrains $J_r\circ r$
when readouts exist, not necessarily the raw readout $r$. For a
partial readout, this distinction also includes its domain.

\begin{proposition}[Exact factorization of a partial readout]
\label{v480:kernel-partial-factorization}
Let $N:A_F\to K$ be total and $f:D\to Y$, $D\subseteq A_F$.
A factor $\bar f:N(D)\to Y$ with $\bar f(N(a))=f(a)$ on $D$
exists exactly when $f$ is constant on the fibres of $N$ inside $D$.
Equality as partial functions, $f=\bar f\circ N$, additionally requires
\[
 D=N^{-1}(N(D)).
\]
If $D=A_F$, this is the total factorization case. The empty-domain
case needs no artificial extension to unobserved classes.
\end{proposition}
\begin{proof}
Define $\bar f(N(a))=f(a)$ for $a\in D$. Fibre constancy is
necessary and sufficient for well-definedness. The composite has
domain $N^{-1}(N(D))$, so equality of partial functions is
equivalent to the displayed domain condition.
\end{proof}

For example, a readout defined only at one of two equivalent
payloads is pairwise compatible wherever defined, yet its quotient
factor composed with $N_F$ becomes defined at both. Exact domain
preservation therefore cannot be inferred from the old pairwise
readout condition alone.

Nor does the proposed display macro establish a sufficient normal
form: take $A_F=\{0,1\}$, closure equivalence equal to equality,
payload extraction $p(j(a))=a$, and $N(a)=0$ for both payloads.
The macro's identity $N(p(j(a)))=N(a)$ holds, but the compatible
readout $f(a)=a$ cannot factor through this constant $N$.
This is a counterexample to sufficiency of the bare macro, not to
the complete admissible-kernel conditions.

\begin{proposition}[A sufficient routed-state condition]
\label{v480:kernel-routed-condition}
Fix an admitted weight record $w$ and an operation
$G_w:A_F^n\to A_F$, $n\geq1$. Suppose $G_w$ respects componentwise
$E_F$-equivalence, and $G_w(a,\ldots,a)E_Fa$. If every input
$a_iE_Fa_0$, then $G_w(a_1,\ldots,a_n)E_Fa_0$. The operation
also descends uniquely to the quotient.
\end{proposition}
\begin{proof}
Congruence compares the given input tuple with $(a_0,\ldots,a_0)$.
The diagonal condition and transitivity give the class-preservation
claim. Congruence is exactly the representative-independence
needed to define the operation on quotient classes.
\end{proof}

Normalized weights alone do not imply this condition: with classes
$\{0,1\}$ and $\{2,3\}$ and weights $(1/2,1/2)$, the constant
operation $G_w(a,b)=2$ is congruent but sends inputs $0,1$ to the
wrong class. By contrast, if the aggregate is specified only at
the level of a vector-valued class readout and all values equal
$v_0$, linear aggregation gives
$\sum_i w_i v_0=v_0$ when $\sum_i w_i=1$. This proves a readout
identity, not the existence or class preservation of a native
state-level aggregator.

\subsection{Shift Compatibility and the Activation Boundary}
\label{v480:kernel-activation-boundary}

For the payload-preserving tag migration
$U(j_0(a))=j_1(a)$, define the shifted closure relation by
transport through $U$. Then $p_1\circ U=p_0$ and both $q_F\circ p_t$
and the finite $N_F\circ p_t$ commute with this display change.
This is a proof for an explicitly declared encoding change, not
an assertion that arbitrary physical frame changes preserve the payload.

The local kernel package is therefore justified when the fragment,
closure equivalence, invariant aspects, readout domains and routed
operations have been supplied with the properties above. The finite
checker validates those supplied finite relations, normal forms,
readout fibres and complete aggregation tables using exact code
equality. It neither infers a native equivalence nor silently
completes an invalid relation.

The original global activation package remains unproved: a native
kernel identification, its aspect/recursive-filter compatibility,
all relevant readouts and routed operations, and a compatible
normal-form construction across the full active surface are still
needed. A quotient class is not a canonical native payload
representative; finite representative selection does not construct
one on an arbitrary proper class. The symbols
\texttt{WARN\_022} and \texttt{WARN\_081} are not declared cleared,
and no external evaluator outcome is inferred from this work.

\paragraph{Explicit positive finite instance.}
Take $A_F=\{0,1,2,3\}$ with classes $\{0,1\}$ and $\{2,3\}$.
Then $N_F=(0,0,2,2)$, and the total class readout
$f=(7,7,11,11)$ factors as $\bar f(0)=7$, $\bar f(2)=11$.
With $c(a)=\lfloor a/2\rfloor$, declare the aspect tuple
$(\mathrm{state},c(a),(\mathrm{class},c(a)))$, constant on each class.
Let $T$ swap $0,1$ and swap $2,3$, and set $G_w(a,b)=a$ for
the admitted weight record $(1/2,1/2)$. The transport stays in
each class, the aggregator is congruent and diagonal-preserving,
and the payload-preserving tag shift commutes with all these
normal-form comparisons. This is a fully specified consistent
local model, not a recovered native QAM witness.
